% week10_compact_gui.tex -- CS2 Week 10 Monday, Oct 19 2026.
% Reasoning Odyssey S04: Compact GUI over a tested model.
% Sources: planning/week-10.md; lessons/week-10-tkinter-model-view-lab.md.
\documentclass[aspectratio=169]{beamer}
\usepackage{../theme/cs2theme}
\usepackage{../theme/cs2commands}

\setcsweek{Week 10}

\title{A Window Is Just a Customer of Your Model}
\subtitle{Compact GUI over a tested model (S04)}
\author{Computer Science II \textbullet\ Dr. Jeremy P. Evert}
\date{October 19, 2026}

\begin{document}

\begin{frame}[plain]
  \titlepage
  \note{Monday of Week 10. The whole class is one idea: the model must work and be provable with no window open. The view is thin. Everything today uses a ten-line model, one label, one button.}
\end{frame}

\begin{destinationframe}{Today's destination}
  \begin{itemize}
    \item A small model class that passes a \texttt{unittest} \emph{without} opening a window.
    \item A Tkinter view: one label, one button, one callback.
    \item The trace: click $\rightarrow$ callback $\rightarrow$ model method $\rightarrow$ new state $\rightarrow$ label updated.
    \item Break the boundary on purpose, then refactor it back.
  \end{itemize}
  \instructorcue{Say the goal out loud: we are not building a desktop app, we are keeping the boundary clean.}
  \note{The plan says larger checkpoint, not a desktop-app engineering unit. Keep the GUI modest. If time runs short, cut the stretch and the 10,000-items discussion, never the trace. Why it matters: a bug you can only see by clicking cannot be tested, and rules that live in the callback die when the window changes. Tie back to the Week 7 contract.}
\end{destinationframe}

\begin{conceptframe}{Core concept: model, view, callback}
  \begin{description}
    \item[Model] Holds the state and the rules. Knows nothing about Tkinter.
    \item[View] Shows the state. Has no rules of its own.
    \item[Callback] The only bridge: receives one event, calls \emph{one} model method, then refreshes the view.
  \end{description}
  \instructorcue{Ask: where does ``you cannot remove an item that is not there'' belong? (Model.)}
  \note{Domain rules never live only in a callback. If a student puts a rule in the callback, it is a candidate for the break-the-boundary slide later.}
\end{conceptframe}

\begin{frame}{Concept + diagram: who knows about whom}
\conceptdiagramkicker
\begin{conceptdiagrambody}
  \begin{concepttext}
    The view \emph{calls} the model.\\[2mm]
    The model never imports Tkinter and never calls the view.\\[2mm]
    That one-way arrow is the whole design.
  \end{concepttext}
  \begin{conceptdiagram}
    \begin{tikzpicture}[node distance=14mm]
      \node[cs2class] (view) {View (Tkinter)};
      \node[cs2class, below=of view] (model) {Model};
      \draw[cs2dataflow] (view) -- node[right, font=\footnotesize]{calls} (model);
    \end{tikzpicture}
  \end{conceptdiagram}
\end{conceptdiagrambody}
\note{Draw it on the board too. One arrow, downward. If a student wants the model to update the label directly, that is the arrow going the wrong way.}
\end{frame}

\begin{frame}[fragile]{Model first: ten lines, no window}
\codeframekicker
\begin{lstlisting}
class Inventory:
    def __init__(self):
        self.items = []

    def restock(self, name):
        self.items.append(name)
        return len(self.items)

    def count(self):
        return len(self.items)
\end{lstlisting}
\instructorcue{Swap in your own world: roster, case file, crew list, shop stock.}
\note{Use whichever world noun the class likes. The method name matters later for the peer swap, so keep restock and count unless students choose a better pair and agree on it.}
\end{frame}

\begin{frame}[fragile]{The test passes with no window}
\codeframekicker
\begin{lstlisting}
import unittest
from inventory import Inventory

class InventoryTest(unittest.TestCase):
    def test_restock_adds_one(self):
        inv = Inventory()
        self.assertEqual(inv.restock("rope"), 1)
        self.assertEqual(inv.count(), 1)

if __name__ == "__main__":
    unittest.main()
\end{lstlisting}
\note{Run it live. Point out: no tkinter import anywhere, and it runs on a server with no display. That is the evidence that the model stands on its own.}
\end{frame}

\begin{frame}[fragile]{Predict before running: the view}
\predictkicker
\begin{lstlisting}
import tkinter as tk
from inventory import Inventory
model, root = Inventory(), tk.Tk()
text = tk.StringVar(value="Items: 0")
tk.Label(root, textvariable=text).pack()
def on_restock():
    n = model.restock("rope")   # route to the model
    text.set(f"Items: {n}")     # then refresh the view
tk.Button(root, text="Restock", command=on_restock).pack()
root.mainloop()
\end{lstlisting}
  \textbf{Predict:} three clicks on Restock. What does the label say?
  \note{Answer: Items: 3. Ask a second question: where is the number stored, the label or the model? (Model. The label only shows it.) Point out the callback has two lines and no decisions. If a student needs an if, ask whether it is a domain rule; if yes, move it into the model and test it.}
\end{frame}

\begin{frame}[fragile]{Trace the event}
\codeframekicker
\begin{lstlisting}
def on_restock():
    print("click -> callback")
    n = model.restock("rope")
    print(f"model.restock -> {n} items")
    text.set(f"Items: {n}")
    print("label updated")
\end{lstlisting}
  \textbf{One click:} click $\rightarrow$ callback $\rightarrow$ model.restock $\rightarrow$ 1 item $\rightarrow$ label updated
  \instructorcue{Students paste their own trace into the World Bible as evidence.}
  \note{This trace is the evidence artifact for today. Remove the prints afterwards if you like, but keep a copy of the output.}
\end{frame}

\begin{misconceptionframe}{Common failure: ``the GUI is the program''}
  \textbf{Symptom:} the only way to know if it works is to run the window and click.

  \vspace{2mm}
  \textbf{Fix:} write the model test first. If you cannot test a rule without a window, the rule is in the wrong place.
  \note{Let students say it before you do. This is the most common Week 10 failure: logic hidden in command=lambda.}
\end{misconceptionframe}

\begin{frame}[fragile]{Break the boundary on purpose}
\begin{comparebody}
  \begin{compareleft}{Boundary broken}
\begin{lstlisting}
items = []   # view owns data
def on_restock():
    items.append("rope")
    n = len(items)
    text.set(f"Items: {n}")
\end{lstlisting}
  \end{compareleft}
  \begin{compareright}{Boundary kept}
\begin{lstlisting}
model = Inventory()
def on_restock():
    n = model.restock("rope")
    text.set(f"Items: {n}")
\end{lstlisting}
  \end{compareright}
\end{comparebody}
\instructorcue{Ask: how would you test the left version without a window?}
\note{Do this live: start from the left version, ask the class what breaks when we want a test, then refactor into the right. The refactor is moving two lines, which is the point: it is cheap if you do it now.}
\end{frame}

\begin{pairdiscussframe}{Think \textbullet\ Pair \textbullet\ Discuss}
  Pick the single most useful action in \emph{your} world: add a colonist, close a case, restock an item. Name the model method, then the one label and one button that would call it.
  \stopandwork{5 minutes. Write the method name and the label text on paper. No Tkinter yet.}
  \note{Circulate. Look for method names that describe a domain action, not a GUI action. add_colonist is good, on_button_click is not.}
\end{pairdiscussframe}

\begin{designdecisionframe}{Peer swap: plug in your partner's model}
  Trade model files with a partner who used the \emph{same method names}. Run your view on their model.
  \begin{itemize}
    \item It works: the contract held (your Week 7 contract test, again).
    \item It fails: the method names or return values differ. Decide which side should change.
  \end{itemize}
  \note{A very small contract test: assert hasattr(model, "restock") and isinstance(model.restock("x"), int). The failure is the lesson, not an accident to avoid.}
\end{designdecisionframe}

\begin{worldbibleframe}{World Bible entry}
  \begin{worldbibleentry}{Rule: The view never holds the rules}
    \wbfield{Assumption}{Every change to world state goes through one model method.}
    \wbfield{Invariant}{The model test passes with no window; the label always shows the model's value.}
    \wbfield{Revision}{Date, and what you moved out of a callback into the model.}
  \end{worldbibleentry}
  \note{Students record one rule they moved, with the date. This is the World Bible evidence for the week, together with the click trace.}
\end{worldbibleframe}

\begin{odysseygateframe}{Odyssey checkpoint note}
  \textbf{This week's checkpoint:} render real model state in a modest view and route one meaningful event through a callback to a model operation. Keep the boundary clear and show the model works on its own.

  \vspace{2mm}
  \textbf{Evidence:} a model class with a window-free \texttt{unittest}; your one-label, one-button view; the click trace in your World Bible.

  \vspace{2mm}
  \small Larger checkpoint, not a desktop-app engineering unit. See the Week 10 gate in Canvas for the rubric; this slide is the pointer, not the rubric.
  \note{Remind them the gate is due Friday 11:59 PM Central. Stretch for fast students: a second event (a text entry that filters the list) or a model undo. Honest limit: with 10,000 items this label-only view is painful, which leads to Week 11 visualization. Do not promise features beyond one label, one button, one callback.}
\end{odysseygateframe}

\begin{recapframe}{Recap}
  \begin{itemize}
    \item Model first, and prove it with a test that needs no window.
    \item The callback routes one event to one model method, then refreshes.
    \item Rules in the view are a bug you cannot test.
    \item Same method names make models swappable (Week 7 contract).
  \end{itemize}
\end{recapframe}

\begin{exitframe}{Exit question}
  \Large\color{cs2primary}
  \textbf{Name one rule in your world that must live in the model, and say how you would test it without opening a window.}

  \vspace{4mm}
  \normalsize\normalfont\color{cs2ink}
  \textbf{Next move:} bring a working model and its test to Wednesday.
\end{exitframe}

\end{document}
